\chapter{White Blood Cell Count (WBC)}

\section*{What Is This Test?}
The white blood cell (WBC) count, also called the leukocyte count, measures the total number of white blood cells (leukocytes) per unit volume of whole blood. White blood cells are the cellular mediators of the immune system and are responsible for defending the body against infection, inflammation, and foreign substances. There are five major types of leukocytes circulating in the peripheral blood: neutrophils (which phagocytose bacteria and fungi), lymphocytes (which mediate adaptive immunity through B cells, T cells, and natural killer cells), monocytes (which differentiate into tissue macrophages and dendritic cells), eosinophils (which combat parasitic infections and participate in allergic responses), and basophils (which release histamine and heparin in inflammatory and allergic reactions). The total WBC count is a core component of the complete blood count and is used as a screening test for infection, inflammation, bone marrow function, and hematologic malignancy. An elevated WBC count (leukocytosis) most commonly reflects an appropriate response to infection or inflammation, but may also indicate leukemia or other myeloproliferative disorders. A decreased WBC count (leukopenia) suggests bone marrow suppression, overwhelming infection, autoimmune disease, or drug effect. The WBC count must always be interpreted in conjunction with the differential count, which provides the percentage and absolute number of each leukocyte type.

\section*{Alternative Names}
WBC; Leukocyte Count; White Count; White Blood Cell Number; Total Leukocyte Count; TLC

\section*{Principle}
The WBC count is measured on automated hematology analyzers using either the electrical impedance (Coulter) principle or optical flow cytometry. In impedance counting, whole blood is diluted with a lysing reagent that eliminates red blood cells while preserving white blood cells. The white cell suspension passes through a small electrical aperture. Each cell generates an electrical resistance pulse proportional to its size. The number of pulses within a calibrated counting volume is counted to yield the WBC count. In optical flow cytometric methods (Sysmex XN, Siemens ADVIA), white blood cells are characterized by laser light scatter and fluorescence after staining with specific dyes. Some analyzers use a fluorescent nucleic acid dye that differentially stains white blood cell nuclei, and the WBC count is obtained from the total number of fluorescent-stained events. The analyzers count 10,000--30,000 white blood cells per sample, giving excellent precision (coefficient of variation 2--3\%). Nucleated red blood cells (nRBCs), if present, are counted as white cells by impedance methods, causing a falsely elevated WBC count. Modern analyzers detect and correct for nRBCs using fluorescence staining or other methods. For manual counting, blood is diluted 1:20 with a WBC diluting fluid (acetic acid or hydrochloric acid to lyse red blood cells), and cells are counted in a Neubauer hemocytometer chamber. Manual counting is now rarely performed, reserved for situations where automated results are unreliable, such as very low WBC counts or cryoglobulinemia.

\section*{History}
White blood cells were first observed by Antonie van Leeuwenhoek in the late 17th century, but systematic study began with Paul Ehrlich in the 1870s--1880s, who used his aniline dye staining techniques to distinguish different types of leukocytes (neutrophils, eosinophils, basophils). Ehrlich also recognized that WBC counts could be elevated in infection and leukemia. The development of the hemocytometer counting chamber (Thoma-Zeiss, 1878) enabled manual WBC counting. For the first half of the 20th century, WBC counts were performed manually in every hematology laboratory, with technologists diluting blood with acetic acid-based WBC diluting fluid and counting cells on the hemocytometer grid. This manual method had a coefficient of variation of approximately 15--20\% when counting 100 cells. The introduction of the Coulter Counter in 1956 for automated WBC counting was transformative, dramatically improving precision, decreasing turnaround time, and allowing high-throughput processing of large numbers of samples. The total WBC count remains one of the most frequently ordered laboratory tests in clinical medicine.

\section*{Why This Test Is Ordered}
\begin{itemize}
  \item Screening for infection --- leukocytosis (WBC >11,000/$\mu$L) with neutrophilia suggests bacterial infection; leukopenia may suggest viral infection or overwhelming sepsis
  \item Diagnosis of leukemia and other hematologic malignancies --- markedly elevated WBC with blasts, abnormal differential, and cytopenias
  \item Monitoring response to chemotherapy and radiation therapy --- expected decline in WBC count (nadir at 7--14 days after chemotherapy) with subsequent recovery
  \item Evaluation of unexplained fever, weight loss, fatigue, night sweats, or recurrent infections
  \item Assessment of bone marrow function --- in pancytopenia (aplastic anemia, myelodysplasia, bone marrow infiltration)
  \item Monitoring immunosuppressive therapy --- tracking absolute neutrophil count to assess infection risk
  \item Preoperative screening --- part of standard preoperative laboratory evaluation
  \item Detection of drug-induced leukopenia --- agranulocytosis from clozapine, antithyroid drugs, sulfonamides, phenytoin, chemotherapy
  \item Screening for congenital immunodeficiency disorders --- severe congenital neutropenia (Kostmann syndrome), cyclic neutropenia
  \item Evaluation of allergic conditions and parasitic infections --- eosinophilia reflected in elevated total WBC
\end{itemize}

\section*{Primary vs Secondary Test}
The WBC count is a primary screening test and core component of the complete blood count, used for the detection and monitoring of infection, inflammation, hematologic malignancy, and bone marrow function.

\section*{How This Test Is Performed}
A venous blood sample is collected in a lavender-top (EDTA) tube from a peripheral vein, usually the antecubital fossa. Approximately 2--5 mL of blood is drawn, and the tube is gently inverted 8--10 times to mix thoroughly with the EDTA anticoagulant. The sample is processed on an automated hematology analyzer, which lyses red blood cells, counts the remaining white blood cells, and reports the total WBC concentration. The automated analysis takes 1--2 minutes to complete. Manual WBC counting (rarely needed) involves diluting 20 $\mu$L of EDTA blood in 380 $\mu$L of WBC diluting fluid (acetic acid or 1\% hydrochloric acid), allowing 5 minutes for red cell lysis, loading the Neubauer hemocytometer, and counting cells on 4 large corner squares (each 1 mm\textsuperscript{2}). The venipuncture procedure takes less than 5 minutes.

\section*{What Preparation Is Needed}
No special preparation is required. Fasting is not necessary for a WBC count. Patients should avoid vigorous exercise for 30 minutes before blood collection, as intense physical activity can cause transient neutrophilia due to demargination of neutrophils from the endothelial wall (shift leukocytosis). The patient should inform the healthcare provider of all medications, with particular attention to corticosteroids (which cause neutrophilia, lymphopenia, and eosinopenia), lithium (which causes neutrophilia), beta-agonists like epinephrine (which demarginate neutrophils), chemotherapy agents, antipsychotics (clozapine, which can cause agranulocytosis), and antithyroid medications. Recent infection, vaccination, or surgery can alter the WBC count and should be disclosed.

\section*{Contraindications and Precautions}
\begin{itemize}
  \item No absolute contraindications. Important pre-analytical factors affecting results include: EDTA is the standard anticoagulant; the sample must be thoroughly mixed to avoid microclots that can trap white blood cells and falsely decrease the count;hemolyzed or clotted specimens are unacceptable; prolonged storage at room temperature (>24 hours) causes white blood cell degeneration and decreased counts; nucleated red blood cells (nRBCs) are counted as white blood cells by impedance analyzers and must be corrected if present --- the corrected WBC = (measured WBC $\times$ 100) / (100 + nRBCs per 100 WBCs); cryoglobulins may cause leukocyte clumping and spurious counts (warming the sample to 37\textdegree{}C resolves this); very high WBC counts (>100,000/$\mu$L) may be outside the linear range of some analyzers and require dilution; lipemia, hyperbilirubinemia, and paraproteinemia may interfere with some automated methods. Samples should be analyzed within 4 hours at room temperature or 24 hours if refrigerated at 2--8\textdegree{}C.
\end{itemize}
\section*{Risks}
Minimal risks associated with venipuncture: mild pain or stinging at the needle insertion site, bruising (ecchymosis) or hematoma at the puncture site, lightheadedness or vasovagal syncope, and extremely rarely, infection or nerve injury.

\section*{What Happens After the Test}
WBC count results are typically available within 30--60 minutes as part of the CBC. The WBC count is interpreted in conjunction with the differential. An elevated or decreased WBC count prompts review of the differential to identify which cell type is responsible. The healthcare provider determines whether the abnormality is reactive (appropriate response to infection, inflammation, or medication) or represents a primary hematologic disorder. Markedly abnormal results (WBC >100,000/$\mu$L with blasts, WBC <500/$\mu$L) are reported as critical values requiring immediate clinical action. Further testing is guided by the pattern of abnormality: cultures and imaging for suspected infection, flow cytometry and bone marrow biopsy for suspected leukemia, and medication review for suspected drug-induced leukopenia.

\section*{Understanding Results}

\subsection*{Below Normal (Leukopenia: WBC <4,500/$\mu$L)}
\begin{itemize}
  \item \textbf{Viral infections:} Influenza, hepatitis, HIV, Epstein-Barr virus, cytomegalovirus, measles, rubella, dengue fever, and COVID-19 commonly cause transient leukopenia with lymphopenia.
  \item \textbf{Overwhelming bacterial sepsis:} Severe bacterial infection can paradoxically consume and deplete neutrophils faster than the bone marrow can produce them, causing leukopenia (a poor prognostic sign).
  \item \textbf{Drug-induced agranulocytosis:} Idiosyncratic, severe neutropenia (ANC <200/$\mu$L) from clozapine, antithyroid drugs (propylthiouracil, methimazole), sulfonamides, phenytoin, carbamazepine, NSAIDs, chloramphenicol, and ticlopidine. Requires immediate drug discontinuation.
  \item \textbf{Chemotherapy and radiation therapy:} Predictable, dose-dependent myelosuppression with nadir at 7--14 days after administration. Filgrastim (G-CSF) may be used to accelerate neutrophil recovery.
  \item \textbf{Aplastic anemia:} Pancytopenia with leukopenia, anemia, and thrombocytopenia from bone marrow failure. Causes include radiation, benzene, chemotherapy, viral hepatitis, autoimmune destruction, and idiopathic.
  \item \textbf{Myelodysplastic syndromes:} Ineffective hematopoiesis with cytopenias, most commonly in older adults.
  \item \textbf{Autoimmune neutropenia:} Isolated neutropenia from anti-neutrophil antibodies, may be primary or associated with systemic lupus erythematosus, rheumatoid arthritis (Felty syndrome), or Sj\"{o}gren syndrome.
  \item \textbf{Hypersplenism:} Splenomegaly from cirrhosis, portal hypertension, Gaucher disease, or lymphoma causes sequestration and premature destruction of circulating leukocytes.
  \item \textbf{Bone marrow infiltration:} Leukemia, lymphoma, multiple myeloma, myelofibrosis, granulomatous disease (tuberculosis, sarcoidosis), or metastatic carcinoma replacing normal marrow.
  \item \textbf{Nutritional deficiencies:} Severe vitamin B12 or folate deficiency (megaloblastic anemia with pancytopenia), copper deficiency, protein-calorie malnutrition, chronic alcoholism.
  \item \textbf{Congenital neutropenia:} Kostmann syndrome (severe congenital neutropenia), cyclic neutropenia (21-day oscillations of neutrophil count), Shwachman-Diamond syndrome.
\end{itemize}

\subsection*{Above Normal (Leukocytosis: WBC >11,000/$\mu$L)}
\begin{itemize}
  \item \textbf{Acute bacterial infection:} The most common cause of leukocytosis. Neutrophilia (absolute neutrophil count >7,000/$\mu$L) with left shift (bands, metamyelocytes, myelocytes), toxic granulation, and D\"{o}hle bodies. Typical infections include pneumonia, pyelonephritis, abscess, meningitis, appendicitis, cholecystitis, and septic arthritis.
  \item \textbf{Acute inflammation and tissue necrosis:} Burns, trauma, surgery, myocardial infarction, pulmonary embolism, pancreatitis, and gout.
  \item \textbf{Physiological leukocytosis:} Vigorous exercise (epinephrine-induced demargination), pregnancy (especially third trimester), labor and delivery, emotional stress, and smoking.
  \item \textbf{Drug-induced leukocytosis:} Corticosteroids (demargination of neutrophils from vessel walls, peaking 4--6 hours after dose), lithium (stimulates granulopoiesis), epinephrine, G-CSF (filgrastim) or GM-CSF (sargramostim) administration.
  \item \textbf{Chronic myeloid leukemia (CML):} Markedly elevated WBC (often >50,000/$\mu$L, may exceed 200,000/$\mu$L) with a full spectrum of myeloid maturation (blasts, promyelocytes, myelocytes, metamyelocytes, bands, segmented neutrophils). Basophilia is a hallmark feature. Philadelphia chromosome t(9;22) and BCR-ABL1 fusion gene confirm the diagnosis.
  \item \textbf{Acute leukemia:} WBC may be elevated, normal, or decreased. When elevated, blast cells constitute the predominant population (lymphoblasts in ALL, myeloblasts in AML). Auer rods are pathognomonic for AML.
  \item \textbf{Chronic lymphocytic leukemia (CLL):} Lymphocytosis with absolute lymphocyte count >5,000/$\mu$L, smudge cells on peripheral smear, monoclonal B cells expressing CD5, CD19, and CD23 by flow cytometry.
  \item \textbf{Myeloproliferative neoplasms:} Polycythemia vera, essential thrombocythemia, and primary myelofibrosis may present with mild to moderate neutrophilia and thrombocytosis.
  \item \textbf{Leukemoid reaction:} Extreme neutrophilia (>50,000/$\mu$L) with left shift mimicking CML, but occurring as a reactive response to severe infection, malignancy, or acute hemolysis. Distinguished from CML by elevated leukocyte alkaline phosphatase (LAP) score and absence of Philadelphia chromosome.
  \item \textbf{Tissue necrosis and malignancy:} Solid tumors (especially lung, gastric, pancreatic, and renal cancers) can produce leukocytosis through paraneoplastic cytokine production (G-CSF).
  \item \textbf{Rebound leukocytosis:} Following recovery from agranulocytosis or chemotherapy-induced myelosuppression.
  \item \textbf{Post-splenectomy:} Loss of splenic sequestration and removal of senescent leukocytes causes mild, persistent leukocytosis and thrombocytosis.
\end{itemize}

\section*{Normal Results}
\begin{itemize}
  \item \textbf{Adult males and females:} 4,500--11,000/$\mu$L (4.5--11.0 $\times$ 10\textsuperscript{9}/L)
  \item \textbf{Newborns (first week):} 9,000--30,000/$\mu$L (physiological leukocytosis at birth, declining over weeks)
  \item \textbf{Infants (1--12 months):} 6,000--17,500/$\mu$L
  \item \textbf{Children (1--6 years):} 5,000--15,000/$\mu$L
  \item \textbf{Children (6--12 years):} 4,500--13,500/$\mu$L
  \item \textbf{Pregnant women (third trimester):} 5,800--13,200/$\mu$L (physiological increase; count may rise further during labor, reaching 25,000/$\mu$L)
  \item \textbf{African American individuals:} May normally have slightly lower WBC counts (lower limit may extend to 3,500/$\mu$L) due to benign ethnic neutropenia
\end{itemize}

\section*{Follow-up Tests}
\begin{itemize}
  \item \textbf{WBC differential (automated and/or manual):} Essential for identifying which cell type is increased or decreased, presence of left shift, blasts, or abnormal cells
  \item \textbf{Peripheral blood smear:} Direct morphological examination for toxic granulation, D\"{o}hle bodies, blast cells, Auer rods, smudge cells, atypical lymphocytes, or hairy cells
  \item \textbf{Blood cultures and infectious disease testing:} When infection is suspected --- bacterial cultures (aerobic and anaerobic), viral serologies (EBV, CMV, HIV, hepatitis panel), and PCR testing
  \item \textbf{Flow cytometry/immunophenotyping:} For suspected leukemia or lymphoma --- determines cell lineage, clonality, and specific diagnostic markers (CD markers)
  \item \textbf{Bone marrow aspiration and biopsy:} When unexplained leukocytosis with blasts, pancytopenia, or suspected myelodysplasia/leukemia is present
  \item \textbf{Cytogenetic analysis and FISH:} Philadelphia chromosome t(9;22), PML-RARA t(15;17), and other genetic abnormalities in hematologic malignancy
  \item \textbf{Molecular testing:} BCR-ABL1 (CML), JAK2 V617F (myeloproliferative neoplasms), mutation panels
  \item \textbf{Leukocyte alkaline phosphatase (LAP) score:} To differentiate CML (low score) from leukemoid reaction (high score)
  \item \textbf{CRP and ESR:} Markers of inflammation when infection or inflammatory disease is suspected
  \item \textbf{Imaging studies:} Chest X-ray, CT scan, ultrasound for suspected infection, abscess, or malignancy
  \item \textbf{Medication review:} Evaluate for drugs known to cause agranulocytosis or leukocytosis
\end{itemize}

\section*{References}
\begin{enumerate}
  \item Tierney LM, Saint S, Drazen JM. CURRENT Medical Diagnosis \& Treatment. McGraw-Hill; 2023.
  \item Wallach J. Interpretation of Diagnostic Tests. 11th ed. Wolters Kluwer; 2022.
  \item Fischbach FT, Dunning MB. A Manual of Laboratory and Diagnostic Tests. 11th ed. Wolters Kluwer; 2023.
  \item Burtis CA, Ashwood ER, Bruns DE. Tietz Textbook of Clinical Chemistry and Molecular Diagnostics. 7th ed. Elsevier; 2023.
  \item Kaushansky K, Lichtman MA, Prchal JT, et al. Williams Hematology. 10th ed. McGraw-Hill; 2021.
  \item Bain BJ, Bates I, Laffan MA. Dacie and Lewis Practical Haematology. 12th ed. Elsevier; 2017.
  \item Bain BJ. Blood Cells: A Practical Guide. 6th ed. Wiley Blackwell; 2022.
  \item Boxer LA. How to approach neutropenia. Hematology Am Soc Hematol Educ Program. 2012;2012:174--182.
\end{enumerate}

\clearpage
\section*{Regulatory Status and Authorization Date}
This chapter describes a test category, not a specific commercial product. FDA, CDSCO, EU IVDR, and MHRA approval or registration dates therefore cannot be assigned without the assay manufacturer, product name, instrument, and intended use. For laboratory-developed tests, an individual agency approval date may not exist. See the Preface for the interpretation of regulatory status.

\clearpage
